\section{Black Holes and Hawking Radiation}
The horizon as the boundary of the Sphaera
\paragraph*{Definition}
10.1
(Schwarzschild horizon in Sphaera coordinates). The Schwarzschild horizon rS = 2GM/c2 corresponds to $\beta$ $\to$1 in the Sphaera metric gtt = 1 -$\beta$2. In sector coordinates: the horizon is where z-$\to$$\infty$.
\paragraph*{Proposition 10.2 (Horizon as cascade singularity). The cascade operator $\hat K$}
is undefined at z-= $\infty$: the Sphaera cannot project beyond its own horizon. The horizon is the boundary of the Sphaera's jurisdiction.
\paragraph*{Proof. $\hat K$ = C$\mu$$\nu$ $\circ$$\Pi$J requires $\|$z-$\|$< $\infty$(finite sector displacement).}
At z-$\to$$\infty$: $\beta$ $\to$1, $\gamma$Lorentz $\to$$\infty$, the Lorentz transformation is singular, and the projection $\Pi$J ceases to be bounded. Hawking radiation as compensator output
\paragraph*{Theorem 10.3 (Hawking radiation from the triolic structure). At the horizon, a virtual matter--tachyon pair from the triolic structure (Theorem 3.2) is}
promoted to a real pair:
- Tachyon ($\beta$ = -$\beta$sym): falls into the black hole.
- Matter ($\beta$ = +$\beta$sym): escapes to infinity.
- Their compensator output C(f) =
2(f + Jf): a photon with $\beta$ = 0, z-= 0 --- the Hawking radiation. The temperature of the Hawking photon in the photon frame (z-= 0) is: TSphaera = $\gamma$string $\cdot$ m2 Pc2 8$\pi$MkB = $\gamma$string $\cdot$ THawking.
\paragraph*{Proof. The Hawking temperature (standard derivation) is}
TH = $\hbar$c3 8$\pi$GMkB = m2 Pc2 8$\pi$MkB . The Sphaera produces the Hawking photon in the photon frame where $\gamma$ is measured as $\gamma$string (not $\gamma$BI, which is the matter-frame value). Hence TSphaera = $\gamma$string $\cdot$ TH. The ratio TSphaera/THawking = $\gamma$string = ln 2/($\pi$ $\sqrt{}$ 3) $\approx$0.1274 is constant across all black hole masses (verified numerically for M = mP and M = M$\odot$).
\paragraph*{Remark (Interpretation). THawking is the temperature measured by a matter}
observer far from the black hole. TSphaera is the temperature in the photon frame (z-= 0). The ratio $\gamma$string is the frame transformation factor between them --- exactly as $\gamma$BI is the factor between the matter and photon frames for other observables (Theorems 9.5 and 9.6). The Hawking radiation is the Sphaera projecting through the horizon: $\hat K$ fails at z-= $\infty$, but the compensator C still produces a well-defined output at z-= 0 from the pair (fmatter, Jftachyon).
\paragraph*{Remark (Natural observer distance from the horizon). Theorem 9.6 places}
the matter observer at $\beta$sym $\approx$0.5493. In Schwarzschild coordinates, this corresponds to r/rS = 1/(1 -$\beta$2 sym) $\approx$1.43: our natural distance from a black hole is 1.43 rS, just outside the horizon. The origin of the tachyon: causal direction
\paragraph*{Theorem 10.4 (Hawking radiation as the voice of the core). The Hawking}
tachyon is not created at the horizon. It is the black hole core --- which has z-$\to$$\infty$and cannot be projected by $\hat K$ --- attempting to reach z-= 0 through the only available channel: tunnelling through the horizon. The two descriptions are CPT-dual and physically identical: Outside view (t forward) Inside view (t reversed) t1: vacuum t3: core at z-= $\infty$ t2: virtual pair at horizon t2: tachyon at horizon t3: tachyon falls in, matter escapes t1: Hawking photon emitted The cascade $\hat K$ has no preferred time direction: it is the equilibrium between z-= 0 (outside) and z-= $\infty$(core). The Hawking pair is the discharge of this tension.
\paragraph*{Structural argument. Outside: z-grows from 0 (far away) to $\infty$(horizon). $\hat K$}
attempts to project z-$\to$0 but fails at the horizon. Tension forces a real pair. Inside: z-= $\infty$(core, singular compression). $\hat K$ attempts z-$\to$0 but the horizon blocks outward projection. The only escape: a core fraction tunnels as a tachyon ($\beta$ = -$\beta$sym) through the horizon. Under J-reflection, this tachyon appears outside as matter ($\beta$ = +$\beta$sym). The compensator C(f) = 1 2(f + Jf) produces the Hawking photon at $\beta$ = 0. CPT: C exchanges matter and tachyon; P exchanges $\beta$ $\leftrightarrow$-$\beta$ (sides of the horizon); T reverses creation and annihilation. Both descriptions are related by CPT.
\paragraph*{Remark (Black hole as inverted Sphaera). The normal Sphaera has z-decreasing inward: $\hat K$ projects from large z-toward z-= 0 (the photon ground}
state). A black hole inverts this: z-increases inward, from z-= 0 (far away) to z-= $\infty$(horizon and core). The cascade $\hat K$ cannot function inside the horizon; the interior is a region where the Sphaera does not exist. The horizon is where the Sphaera is maximally stressed. Hawking radiation is the Sphaera discharging at its boundary. The black hole evaporates because the core is the Sphaera --- and the Sphaera always tends toward z-= 0. The tachyon is not a particle falling in. It is the core, speaking through the only open channel. Bekenstein--Hawking entropy
\paragraph*{Remark (The entropy factor). The Bekenstein--Hawking entropy SBH}
= A/(4$\ell$2 P) and the LQG entropy SLQG = A ln 2/(8$\pi$$\gamma$BI$\ell$2 P) agree when $\gamma$BI = ln 2/(2$\pi$) $\approx$0.110. The measured value $\gamma$BI $\approx$0.2375 differs by a factor $\approx$2.15. Within the present framework: the LQG entropy formula is derived in the matter frame ($\gamma$BI), while the Bekenstein--Hawking formula is a photon-frame statement. The discrepancy factor 2.15 should be related to the frame transformation, but its precise form --- unlike the analogous Hawking temperature factor $\gamma$string --- is not yet derived. This is an open quantitative problem.
\paragraph*{Theorem 10.5 (Dual origin of $\gamma$BI). The splitting factor $\gamma$L = $\gamma$BI/$\gamma$string $\approx$}
1.864 has a dual origin [\emph{note: value corrected from 1.197 in v2.29.1}]: it is approximately $1/\bsym$ from the Hawking temperature argument:
$\gamma$L = |{z} triolic sector counting $\cdot$ $\beta$sym |{z} observer geometry $\approx$0.667 $\cdot$ 1.820 = 1.214, within 1.4\% of the measured value $\gamma$L = 1.197. Factor A (2/3, structural): In the triolic structure, three sectors exist (K1, K2, K3). Under Wick rotation, only two are stable ($\gamma$ > 0); the tachyon sector ($\gamma$3 < 0) is projected away by $\hat K$. The fraction of stable sectors is 2/3. $\gamma$string counts all three; $\gamma$BI counts only the two stable ones. Factor B (1/$\beta$sym, geometric): A matter observer at r/rS = 1/(1-$\beta$2 sym) $\approx$
1.43 measures an effective Hawking temperature T(r) = TH/$\beta$sym. Via S =
E/T, this shifts the effective entropy count by $\beta$sym, giving $\gamma$BI = $\gamma$string/$\beta$sym $\approx$
0.1274/0.549 $\approx$0.232 (within 2.4\% of 0.2375).
Both factors are active simultaneously. They are not two competing explanations: they are two aspects of the same collapse from z-= $\infty$to z-= 0, seen from the structural side (A) and the geometric side (B).
\paragraph*{Proof. Factor A. The triolic balance $\gamma$1 + $\gamma$2 + $\gamma$3 = 0 with $\gamma$3 < 0 (Theorem 9.3) means that $\gamma$string = $\gamma$1 is defined over all three sectors. The LQG}
computation counts only spin-network microstates, which correspond to the two stable sectors (bosonic C-sector and fermionic A-sector). The weight 2/3 is the fraction of stable sectors. Factor B. The matter observer at $\beta$sym has r/rS = 1/(1 -$\beta$2 sym). The local Hawking temperature is T(r) = TH(1 -rS/r)-1/2 = TH/$\beta$sym. Since $\gamma$BI is fixed by requiring S = A/(4$\ell$2 P) for the observed temperature, the effective parameter shifts: $\gamma$BI = $\gamma$string $\cdot$ (TH/T(r))-1 = $\gamma$string/$\beta$sym. Combined. $\gamma$L = (2/3) $\cdot$ (1/$\beta$sym) $\approx$1.214 vs measured 1.197: within 1.4\%. The residual 1.4\% is the remaining quantitative open problem.
\paragraph*{Remark (What this resolves). The question "are natural constants phaseshifted?" has a precise answer:}
- $\hbar$, G, c, kB are Lorentz-invariant --- they do not shift.
- $\gamma$BI is a dimensionless ratio, not a fundamental constant. It shifts due to
(A) triolic sector counting and (B) observer geometry.
- The shift is not a frame transformation of the constants; it is a consequence of which states the matter observer has access to, and where the
observer sits relative to the horizon. Nature does not choose between (A) and (B). They are the same statement: a matter observer at $\beta$sym naturally has access to 2 of 3 stable triolic sectors, at a distance that shifts the effective temperature by 1/$\beta$sym.
\paragraph*{Theorem 10.6 (ART--QFT consistency of the Sphaera). The Sphaera biquaternion encodes both General Relativity and Quantum Field Theory, connected by Wick rotation:}
it $\cdot$ 1 + ziei | {z } Minkowski (ART) t7$\to$i$\tau$E ----$\to$$\tau$E $\cdot$ 1 + ziei | {z } Euclidean (QFT) . At the natural observer position r0/rS = 1/(1-$\beta$2 sym), the Schwarzschild metric gives gtt(r0) = -$\beta$2 sym and $\sqrt{}$-g
r0 = r2 0 sin $\theta$ (invariant): no additional geometric factor is hidden in the ART side. The QFT side gives local redshifted quantities: Elocal = E$\infty$/$\beta$sym, Tlocal = TH/$\beta$sym. These are exactly Factors A and B of Theorem 10.5. The ART--QFT conversion is internally consistent.
\paragraph*{Remark (The 1.4\% residual and discreteness). The ART--QFT computation}
confirms the structure of Theorem 10.5 but does not close the 1.4\% gap between (2/3)/$\beta$sym $\approx$1.214 and the measured $\gamma$L $\approx$1.197. The gap is not a one-loop
quantum gravity correction (suppressed by ($\ell$P/r0)2 $\to$0) and not a hidden metric factor ($\sqrt{}$-g is invariant). The most likely source: the transition from the discrete prime lattice $\Gamma$P (Stone--von Neumann, QFT) to the continuous spacetime manifold (ART). This is the continuum limit of lattice quantum gravity --- a known open problem. The 1.4\% is the quantitative signature of this discreteness correction, and the framework cannot close it without a more detailed treatment of the lattice-to-continuum transition.

